% =============================================================================
% Section 12 — Complexity Analysis
% =============================================================================
\section{Complexity Analysis}
\label{sec:complexity}

This section analyses the computational and communication complexity of
\pebble{} as a function of the cluster size $N$ (number of nodes), the total
number of distinct keys $K$, and the total number of mutations $M$ in the
cluster history.

\subsection{Per-Node State Complexity}

\begin{table}[H]
\centering
\caption{In-memory state sizes per node.}
\label{tab:state-complexity}
\begin{tabular}{@{}lll@{}}
\toprule
\textbf{Structure} & \textbf{Size} & \textbf{Notes} \\
\midrule
Current state map ($S$)    & $\Oof{K}$   & One winning mutation per key,
                                            including tombstones \\
Mutation history ($H$)     & $\Oof{M}$   & All mutations ever received;
                                            no GC in v0.1 \\
Version vector ($V$)       & $\Oof{N}$   & One entry per known node \\
\bottomrule
\end{tabular}
\end{table}

The total in-memory footprint per node is $\Oof{K + M + N}$. In practice
$M \gg K$ after steady-state operation (many mutations per key), so the
dominant term is $\Oof{M}$.

\begin{remark}[History growth]
Because \pebble{} v0.1 retains the full mutation history without compaction,
memory usage grows without bound as long as mutations are generated. This is
the primary scalability concern in the current version; see
\cref{sec:limitations} and \cref{sec:future} for the planned GC mechanism.
\end{remark}

\subsection{Version Vector and Missing-Range Computation}

\begin{proposition}
Computing $\mathrm{missing}(V_A, V_B)$ requires $\Oof{N}$ time and produces
at most $N$ range entries.
\end{proposition}

\begin{proof}
The computation iterates over $\mathrm{dom}(V_A)$, which has at most $N$
entries (one per node). For each entry, a constant-time comparison and
conditional append are performed. The result has at most $N$ elements
(one per originating node). \qed
\end{proof}

\subsection{Gossip Round Complexity}

\begin{table}[H]
\centering
\caption{Per-gossip-round message complexity.}
\label{tab:round-complexity}
\begin{tabular}{@{}lll@{}}
\toprule
\textbf{Phase} & \textbf{Messages} & \textbf{Max Payload} \\
\midrule
DIGEST exchange      & 2 messages & $\Oof{N}$ bytes each \\
REQUEST (initiator)  & 1 message  & $\Oof{N}$ ranges \\
MUTATIONS (from B)   & 1 message  & $\Oof{M_{\max}}$ mutations \\
REQUEST (responder)  & 1 message  & $\Oof{N}$ ranges \\
MUTATIONS (from A)   & 1 message  & $\Oof{M_{\max}}$ mutations \\
ACK $\times 2$       & 2 messages & $\Oof{1}$ each \\
\midrule
\textbf{Total}       & \textbf{8 messages} &
  $\Oof{N + M_{\max}}$ bytes \\
\bottomrule
\end{tabular}
\end{table}

where $M_{\max} = \texttt{limits.maxMutationsPerMessage}$ (default: 500).

\paragraph{Worst-case bandwidth.}
If a node is starting fresh with no history and connects to a node with $M$
total mutations, it will require $\lceil M / M_{\max} \rceil$ gossip rounds to
catch up, each transferring $\Oof{M_{\max}}$ mutations. Total catch-up bandwidth:
$\Oof{M}$.

\paragraph{Incremental (steady-state) bandwidth.}
During steady-state operation with $\Delta M$ new mutations since the last sync
with a given peer, a single gossip round transfers $\Oof{\Delta M}$ mutations.
This is the common case.

\subsection{Cluster-Wide Convergence Time}

\pebble{} uses random gossip fanout of 1 per node per interval. This is
equivalent to the \emph{push-pull} gossip model studied by Demers et
al.~\cite{demers1987epidemic}.

\begin{theorem}[Expected convergence time]
\label{thm:convergence-time}
Starting from a single node holding a mutation, the expected number of gossip
rounds for the mutation to reach all $N$ nodes is $\Theta(N \log N)$, assuming
each node selects a uniformly random healthy peer in each round.
\end{theorem}

\begin{proof}[Sketch]
This follows from the standard coupon-collector analysis of epidemic
broadcasting. In round $r$, let $k$ be the number of informed nodes. Each of
the $N - k$ uninformed nodes is contacted by an informed node with probability
$k/N$ per round. By linearity of expectation and the harmonic series, the
expected total rounds until all nodes are informed is:
\[
  \mathbb{E}[\text{rounds}] = N \sum_{k=1}^{N} \frac{1}{k} = \Theta(N \log N).
\]
\qed
\end{proof}

\begin{table}[H]
\centering
\caption{Expected convergence rounds for representative cluster sizes
with \texttt{gossipIntervalMs} = 1{,}000\,ms.}
\label{tab:convergence-table}
\begin{tabular}{@{}rrrr@{}}
\toprule
$N$ & $\mathbb{E}[\text{rounds}]$ & \textbf{Approx.\ wall time} & \textbf{Note} \\
\midrule
3   & $\approx 5.5$ & $\approx 5.5\,\text{s}$  & Small dev cluster \\
5   & $\approx 11$  & $\approx 11\,\text{s}$   & Small production \\
10  & $\approx 29$  & $\approx 29\,\text{s}$   & Medium cluster \\
20  & $\approx 72$  & $\approx 72\,\text{s}$   & Large cluster \\
50  & $\approx 225$ & $\approx 225\,\text{s}$  & Use higher fanout \\
100 & $\approx 519$ & $\approx 519\,\text{s}$  & Use higher fanout \\
\bottomrule
\end{tabular}
\end{table}

\begin{remark}[Fanout scaling]
For clusters larger than $\sim$20 nodes, a gossip fanout greater than 1
(selecting $f > 1$ peers per round) reduces convergence time to
$\Theta(\log N / \log f)$ rounds. This is a planned improvement
(\cref{sec:future}).
\end{remark}

\subsection{Gossip Session Load per Node}

Each node initiates exactly one outbound gossip session per
\texttt{gossipIntervalMs} interval and can receive an unbounded number of
inbound sessions (one per active peer). In a cluster of $N$ nodes:

\begin{itemize}[leftmargin=1.5em]
  \item \textbf{Outbound load}: $\Oof{1}$ session per node per interval (constant).
  \item \textbf{Inbound load}: $\Oof{N}$ sessions per interval, since every
        other node may choose this node as its target.
\end{itemize}

Each inbound session occupies a dedicated TCP connection for the duration of
the 4-step protocol. Because gossip sessions are short-lived and sequential per
initiator, the server must handle up to $N-1$ concurrent inbound connections.
The Node.js event loop can sustain hundreds of concurrent lightweight connections
without saturation.

\subsection{Handshake Complexity}

The 3-way handshake (\cref{sec:nodeauth}) requires 3 round trips and 4
Ed25519 operations (2 signs, 2 verifies) per connection. All operations are
$\Oof{1}$ in node count and dominate only the connection-establishment phase.

\subsection{Snapshot and Recovery Complexity}

\begin{table}[H]
\centering
\caption{Persistence operation complexity.}
\label{tab:persistence-complexity}
\begin{tabular}{@{}ll@{}}
\toprule
\textbf{Operation} & \textbf{Complexity} \\
\midrule
Snapshot write     & $\Oof{K}$ — serialise current state (one entry per key) \\
Snapshot read      & $\Oof{K}$ — deserialise and validate \\
Log append         & $\Oof{1}$ amortised — single NDJSON line + fsync \\
Log replay         & $\Oof{M}$ — replay all entries (bounded by $M$) \\
Recovery (total)   & $\Oof{K + M}$ — snapshot load + log replay \\
\bottomrule
\end{tabular}
\end{table}
